% Generated draft: requires mathematical and source review.
\documentclass[11pt]{article}
\usepackage[margin=25mm]{geometry}
\usepackage{fontspec}
\setmainfont{Latin Modern Roman}
\usepackage{amsmath,amssymb,mathtools,unicode-math}
\setmathfont{Latin Modern Math}
\usepackage{enumitem,needspace,xurl,hyperref}
\setlength{\parindent}{0pt}
\setlength{\parskip}{5pt}
\setlength{\emergencystretch}{3em}
\newcommand{\N}{\mathbb{N}}
\newcommand{\Z}{\mathbb{Z}}
\newcommand{\Q}{\mathbb{Q}}
\newcommand{\R}{\mathbb{R}}
\newcommand{\C}{\mathbb{C}}
\newcommand{\F}{\mathbb{F}}
\newcommand{\A}{\mathbb{A}}
\newcommand{\PP}{\mathbb{P}}
\newcommand{\E}{\mathbb{E}}
\newcommand{\eps}{\varepsilon}
\newcommand{\dd}{\,\mathrm{d}}
\begin{document}
\section*{UnsolvedMath ID: 701374}
\textbf{Source problem code:} catalogue entry
\subsection{Definitions and mathematical statement}
Write \(1\) for the identity element or trivial subgroup, as appropriate. Let \(H=A\rtimes B\) be a Frobenius group with nontrivial kernel \(A\) and nontrivial complement \(B\). Thus the conjugation action satisfies
\[
C_A(b)=\{1\}\qquad(b\in B\setminus\{1\}).
\]
Let \(G\) be a finite nontrivial group, and let \(\alpha:H\to\operatorname{Aut}(G)\) be an action by automorphisms. Interpret the source's notation \(GA\) as the semidirect product \(G\rtimes A\) determined by the restriction of \(\alpha\). Assume that this is a Frobenius group with kernel \(G\) and complement \(A\); equivalently,
\[
\{g\in G:\alpha(a)(g)=g\}=\{1\}
\qquad(a\in A\setminus\{1\}).
\]
These hypotheses make \(A\) and \(B\) finite: the respective actions embed \(A\) into \(\operatorname{Aut}(G)\) and \(B\) into \(\operatorname{Aut}(A)\). No faithfulness of the whole action \(\alpha\) is additionally assumed. Define
\[
C_G(B)=\{g\in G:\alpha(b)(g)=g\text{ for every }b\in B\}.
\]
For a finite group \(X\), let \(\exp(X)\) be the least positive integer \(e\) such that \(x^e=1\) for every \(x\in X\), with \(\exp(\{1\})=1\).

\textbf{Part (b), the remaining question in the supplied record.} Does there exist a function
\[
f:\Z_{\geq1}\times\Z_{\geq1}\longrightarrow\Z_{\geq1}
\]
such that, for every choice of \(A,B,G,\alpha\) satisfying these hypotheses,
\[
\exp(G)\leq f\bigl(|B|,\exp(C_G(B))\bigr)?
\]
The bound must depend only on the two displayed parameters.

\textbf{Part (a), retained as an archived question.} For a nilpotent group \(X\), let \(\operatorname{cl}(X)\) denote its nilpotency class, with \(\operatorname{cl}(\{1\})=0\). Under the same hypotheses, the archived question asks whether the nilpotency class of \(G\) is bounded solely in terms of \(|B|\) and \(\operatorname{cl}(C_G(B))\). The supplied archive reports an affirmative answer to part (a), attributed to N. Yu. Makarenko and P. Shumyatsky. This is a supplied historical assessment, not an independently verified conclusion of this draft.
\subsection{Short English statement}
V. D. Mazurov asks whether a finite group \(G\), acted on by a Frobenius group \(A\rtimes B\) with \(A\) acting fixed-point-freely, has exponent bounded solely by \(|B|\) and the exponent of the subgroup fixed by \(B\). The supplied record retains this as part (b) and reports that the analogous nilpotency-class question, part (a), was answered affirmatively.
\subsection{Sources}
\begin{itemize}
\item Local UnsolvedMath registry: supplied complete record with permanent ID \(701374\), exact source code \texttt{catalogue entry}, and title \emph{Kourovka Notebook Problem 17.72}. Supplied upstream revision: \texttt{d371ee917ac9647fbb06cd6fc529deec6f3dda82}.
\item \emph{The Kourovka Notebook (Unsolved Problems in Group Theory)}, problem 17.72, proposed by V. D. Mazurov; 17th issue (2010), with the supplied record identifying an update in the 21st issue (2026), \texttt{arXiv:1401.0300v48}. The supplied collection description identifies E. I. Khukhro and V. D. Mazurov as editors.
\item Supplied Archive of Solved Problems entry for part (a): affirmative answer attributed to N. Yu. Makarenko and P. Shumyatsky, \emph{Proc. Amer. Math. Soc.} \textbf{138} (2010), 3425--3436. No article title is supplied.
\item The supplied registry and status review label the record \texttt{partially\_solved}, with part (a) archived as solved and part (b) retained in the main section. The status review is dated \texttt{2026-10-16} and records evidence as \texttt{published}. This is supplied dated curation, not an independently verified literature assessment.
\end{itemize}
\end{document}
